\section{Conclusion}
\label{sec:conclusion}

Retrieval inside a merge-based column store is served today by structures that
are built per part and consulted per part, so the cost of a similarity query
follows the number of parts rather than the number of rows that could answer
it. The design here adds two things to that layout: a weighted inverted file
over a basis whose elements may be words, entities, code symbols, learned
concepts, relations or motifs, with block maxima on the store's own granules;
and a summary per region and per part --- centroid, covering radius, per-element
maxima, counting statistics --- aggregated into a structure the planner reads
to exclude parts before it opens them.

The components are mostly not new, and Table~\ref{tab:status} says which are
which. The arrangement is the contribution: bounds from metric indexing and
dynamic pruning, statistics from distributed resource selection, and a
maintenance path through the merges a column store already performs.

What the design owes is a measurement, and the shape of the claim decides the
measurement. Expected cost scaling with the relevant neighbourhood is a
statement about a query distribution at a fixed recall target, so the protocol
holds recall, fits an exponent with an interval to the parts, granules and
bytes touched per correct answer, compares against the granules that hold the
answer, and runs the queries the design is expected to fail on as their own
stratum. Until that lane reports, this paper is a design and a protocol, and
its results section holds placeholders.

\section*{Reproducing}

The lane that fills this paper's tables is in \texttt{hanzoai/benchmarks} and
has not been published. Its first pass runs upstream ClickHouse~26.4.2.10 at
$N = 10^5$ to $10^7$. The harness writes \texttt{results-data.tex} directly; the
tables in \S\ref{sec:results} read it, and a key it does not define prints
\ph{}.
